\subsection{Inteligencia artificial en la enseñanza de las matemáticas}

\subsubsection{Emergencia de la IA en contextos educativos}

En los últimos años, el campo educativo ha incorporado tecnologías emergentes como respuesta a las limitaciones evidenciadas en la enseñanza tradicional. Entre estas innovaciones tecnológicas, la Inteligencia Artificial ha emergido como una herramienta con alto potencial transformador para los procesos de aprendizaje matemático \parencite{panquebanInteligenciaArtificialEducacion2024}. Su capacidad para favorecer experiencias más adaptativas, personalizadas y orientadas al desarrollo del pensamiento matemático avanzado representa un cambio paradigmático en la concepción misma del proceso educativo.

La IA permite generar análisis detallados sobre el desempeño estudiantil, identificar patrones de error específicos y proporcionar retroalimentación inmediata y contextualizada. Estos aspectos resultan particularmente relevantes para áreas de alta abstracción conceptual como la trigonometría \parencite{mirelesAplicacionInteligenciaArtificial2025,hariyantoArtificialIntelligenceAdaptive2025}. A diferencia de las aproximaciones pedagógicas tradicionales, donde el docente debe atender simultáneamente a grupos numerosos de estudiantes con necesidades diversas, los sistemas basados en IA pueden ofrecer acompañamiento individualizado a escala \parencite{suarezDesarrolloHabilidadesCon2023}.

\subsubsection{Evidencia sobre mejora del rendimiento académico}

Diversas investigaciones reportan que los sistemas basados en IA pueden mejorar de manera significativa el rendimiento académico, la motivación intrínseca y la retención de contenidos matemáticos \parencite{calderonmaldonadoInteligenciaArtificialAprendizaje2025,EfectoChatGPTRendimiento2025}. Estas mejoras se logran mediante el ajuste dinámico de las actividades a las necesidades específicas de cada aprendiz y el adaptarse al aprendizaje verdaderamente personalizados \parencite{muirraguiirrazabalImpactoUsoInteligencia2025}.


Aunque una parte considerable de estos estudios se ha desarrollado en entornos universitarios, su relevancia para la educación media es creciente \parencite{parrasanchezPersonalizacionRecursosPara2023}. Las dinámicas de personalización que permiten estos sistemas resultan especialmente valiosas para atender las brechas conceptuales previamente descritas en el aprendizaje de la trigonometría. La capacidad de adaptar el ritmo, la profundidad y la modalidad de presentación de los contenidos según las características individuales de cada estudiante representa una ventaja sustancial sobre los métodos tradicionales de instrucción masiva \parencite{panquebanInteligenciaArtificialEducacion2024}. 

Investigaciones recientes sugieren que estos beneficios trascienden a la mejora en calificaciones. Los estudiantes que interactúan con sistemas adaptativos basados en IA desarrollan mayor autonomía en su aprendizaje, mejoran su percepción de autoeficacia matemática y muestran actitudes más positivas hacia el estudio de las matemáticas \parencite{kuoDevelopingAILearning2026,yaseenImpactAdaptiveLearning2025}.


\subsubsection{Conexión con aplicaciones reales}

Más allá de su aporte a la personalización del aprendizaje, la IA también facilita el establecimiento de conexiones significativas entre los contenidos matemáticos abstractos y aplicaciones concretas en ámbitos tecnológicos contemporáneos \parencite{caoElucidatingSTEMConcepts2023}. Áreas como la robótica, el procesamiento de datos, la navegación autónoma o el análisis de señales dependen críticamente de conceptos trigonométricos que los estudiantes frecuentemente perciben como abstractos y carentes de utilidad práctica.

Esta capacidad de contextualización abre la posibilidad de presentar la trigonometría no solamente como un conjunto de procedimientos algorítmicos desconectados de la realidad, sino como un lenguaje matemático fundamental para comprender y modelar fenómenos del entorno tecnológico y natural. Los sistemas basados en IA pueden simular estos contextos aplicados, permitiendo que los estudiantes exploren las implicaciones prácticas de los conceptos trigonométricos en entornos interactivos y motivadores.

Por ejemplo, un sistema de IA podría simular el diseño de trayectorias para drones, donde los estudiantes deben aplicar funciones trigonométricas para calcular ángulos de ascenso, distancias de desplazamiento y patrones de vuelo periódicos. Esta aproximación no solo incrementa la motivación estudiantil al demostrar la relevancia práctica del conocimiento matemático, sino que también facilita una comprensión más profunda al situar los conceptos abstractos en contextos significativos \parencite{yeungIntegratingDroneTechnology2025}.

